\chapter{Calibration storage and battery behavior}
\section{Persistent fits versus volatile reference data}
The RP2350 build reserves two flash erase sectors, 8192 bytes total, outside the
linked firmware region. Each record stores schema, generation, board identifier,
sensor-profile version, validity mask, gyro count, IMU fits/uncertainties,
magnetic offset/scale and voltage/current/lidar fits.

Explicit little-endian integer/binary32/binary64 encoding avoids struct padding.
A CRC32 checks the data and a separate final page commits it.
Ground barometer reference, live SOC and PID integrals are not persisted.
Ground pressure must be measured again for the current session.

\section{The journal write sequence}
\begin{enumerate}
\item Find the newest complete record with matching CRC and commit page.
\item Prepare the next record in RAM, keeping the previous sector untouched.
\item Erase the other sector and program its data pages.
\item Program the commit page last; then reread the data.
\item Boot ignores incomplete/bad-CRC writes and selects the newest complete record.
\end{enumerate}
This addresses ordinary interrupted writes if the other sector remains intact;
it is not a physical guarantee for every brownout or flash failure. Hardware
power-loss/endurance verification remains necessary. CRC is not authentication.

Board/profile mismatch or invalid fit ranges require recalibration.
A sensor replacement on the same board must invalidate the fit/profile explicitly;
board identity alone cannot identify every changed component.
Clearing calibration writes a newer empty record so an older fit is not silently
resurrected on normal reboot.

\section{Verify the gyro after boot}
Saved bias is loaded, then fresh stationary data verify it.
For old/new bias fits and their estimated variances:
\begin{equation}
|\hat b_{\rm new}-\hat b_{\rm old}|>
\max(.05,\ 5\sqrt{P_{\rm old}+P_{\rm new}})
\end{equation}
in degrees/s requests a refresh on any axis. Missing gyro fits also refresh.
The current criterion is provisional engineering policy, not a measured
temperature model. Unstable captures do not complete verification.

Only after verification and required fits exist can the sample helper mark an
IMU sample calibrated. Missing accelerometer scale needs six-face calibration;
level-only capture cannot recover all scale factors. The Pico's live sensor
driver is still absent, so this code never invents capture data.

\section{Programming conditions and commands}
Flash programming conflicts with execute-in-place reads. The RAM write callback
uses Pico SDK \texttt{flash\_safe\_execute}; the other core initializes participation
in lockout. Writes are accepted only on core 0 while the caller is disarmed.
Linker exclusion and a runtime overlap check protect the firmware area.
Unchanged fits do not erase again; failed automatic saves retry no faster than
every five seconds. After three automatic failures, retries stop until a manual
S command or a newly installed fit resets the retry allowance.

\begin{center}\begin{tabular}{@{}lp{.75\linewidth}@{}}\toprule
USB command & Meaning\\\midrule
C & Print loaded/missing calibration, gyro verification and unsaved status.\\
R & Persist an empty record and request recapture/recalculation.\\
S & Save valid measured fits, if any exist.\\\bottomrule
\end{tabular}\end{center}
These commands do not generate observations. Full-chip erase loses calibration.
Configure actual physical flash size and board before flashing a custom PCB.

\section{Requested SOC policy}
For charge fraction $z$, the selected policy is
\begin{equation}
A(z)=\begin{cases}
\text{normal},&100z>35,\\
\text{warning},&20<100z\le35,\\
\text{latched disarm},&100z\le20.
\end{cases}
\end{equation}
The 35-percent warning flag is available for telemetry; current firmware prints
transitions on USB when a monitor supplies data. Handheld LoRa notification is
not yet connected. There is no live voltage/current acquisition, so the firmware
does not pretend it knows real pack SOC.

At 20 percent the C mathematical demands are inhibited, regardless of height.
The latch stays set if a later SOC estimate rebounds. An explicit reset while
disarmed starts a new battery/session policy state. Critical external BMS/cell
alerts also latch disarm; ordinary voltage/current warnings notify.
There is no automatic battery landing request. Physical arming/ESC output remains
disabled.

The simulation bridge accepts a manually supplied estimated SOC percentage
(Float32) on this topic:
\begin{quote}\small\texttt{/qcopter/battery\_soc\_percent}\end{quote}
It publishes UInt8 flags on
\texttt{/qcopter/battery\_policy}: bit 0 warning, bit 1 latched disarm.
These reach the same host C instance that computes the flight demands.
No new ROS flight has exercised this interface yet; it does not manufacture
a battery discharge measurement.

\chapter{ROS, Gazebo and implementation boundaries}
\section{Current simulation feedback path}
\begin{center}
\begin{tikzpicture}[>=Latex,node distance=7mm,
box/.style={draw=ink,rounded corners,fill=soft,text width=9cm,align=center,minimum height=9mm}]
\node[box] (g){Native Gazebo state, physics and RGB};
\node[box,below=of g] (s){Sensor model and shared C++ mission core};
\node[box,below=of s] (d){Local IPC and ROS 2 DDS adapters};
\node[box,below=of d] (c){Host-compiled C filters, six controllers and mixer};
\node[box,below=of c] (m){Motor dynamics, thrust/torque, native link wrench};
\draw[->] (g)--(s);\draw[->] (s)--(d);\draw[->] (d)--(c);\draw[->] (c)--(m);
\draw[->] (m.east)--++(1,0)|-(g.east);
\end{tikzpicture}
\end{center}
The host library runs production C math, not an RP2350 emulator.
Virtual motor gating is separate from physical arming.
The recorded flight uses assumed actuator dynamics, synthetic inertial/height
observations, native camera pixels and simulated local odometry. Its downward
range is an analytical floor ray. The full standard ROS C++ graph was built and
checked separately, not silently substituted as the recorded flight path.

\section{Future physical companion}
\begin{equation*}
\text{ROS mission}\to\text{USB/UART bridge}\to\text{RP2350 parser}
\to\text{local 500 Hz control}.
\end{equation*}
Physical framing, sequence/session, CRC, freshness, status/acknowledgement,
clock handling and command authority remain work. ROS angle units are radians;
the current C interface uses degrees. Convert at one documented boundary.
A slower companion setpoint stream does not clock every inner PID iteration.

\section{Build evidence and limits}
The current simulator launcher is
\begin{quote}\ttfamily\small
powershell -ExecutionPolicy Bypass -File tools/run\_native\_garden.ps1
\end{quote}
After editing the C source, rebuild host libraries using
\texttt{tools/build\_sitl.ps1}. The separate Pico build produces ELF/UF2.
No physical board was flashed. Compilation does not prove flash durability,
sensor acquisition, telemetry delivery or actual ESC inhibition. No new flight
or implementation test suite was executed for this change.

\chapter*{References and reproducibility}
\addcontentsline{toc}{chapter}{References and reproducibility}
\begin{itemize}
\item Raspberry Pi Pico SDK flash APIs and execution safety:
\url{https://www.raspberrypi.com/documentation/pico-sdk/hardware.html}.
The installed \texttt{pico/flash.h} and linker/flash implementation were inspected.
\item NIST uncertainty evaluation:
\href{https://www.nist.gov/pml/nist-technical-note-1297}{Technical Note 1297},
\href{https://physics.nist.gov/cuu/Uncertainty/typea.html}{Type A evaluation},
\href{https://physics.nist.gov/cuu/Uncertainty/typeb.html}{Type B evaluation}.
\item ROS axes/units:
\href{https://raw.githubusercontent.com/ros-infrastructure/rep/master/rep-0103.rst}{REP 103: standard units and coordinate conventions}.
\item TDK InvenSense MPU6000/MPU6050 specification: units, noise and temperature
dependence; actual sensor configuration still needs characterization.
\item Reproducible local model calculations:
\texttt{tools/calculate\_current\_control\_design.py} and the corresponding
\texttt{artifacts/model-design/current/calculations.json}.
\item Single-file runtime source:
\texttt{flight\_controller/flight\_controller.c}.
\end{itemize}
\important{Teaching numbers and scenario coefficients are labelled assumptions.
No synthetic example is presented as measured aircraft data. Preserve raw
measurements, reference uncertainty and source/configuration with each result.}
